{\color{red}
\section{Pole table and account of all $1/\epsilon$ poles (P3)}\label{sec:dv-table}
\noindent[\textbf{New section; fix P3} (fix\_poles\_dglap Steps 3--7; exact rational arithmetic in \texttt{fix/poles/s1\_pole\_algebra.py}, regions in \texttt{s2\_uv\_ir\_regions.py}).] \textbf{Regularisation:} transverse dimension $d=2-\epsilon$ for the loop and the collinear integrals. A scaleless integral is $\int d^{d}p/p^{2}=\pi\big(\frac{2}{\epsilon_{\rm UV}}-\frac{2}{\epsilon_{\rm IR}}\big)$; it is zero in the author's single-$\epsilon$ bookkeeping ($\epsilon_{\rm UV}=\epsilon_{\rm IR}$). Both bookkeepings are given; they lead to the same finite result.

\subsection*{Step 3: the normalisation of the dressed measured gluon (not written in the rows)}
\noindent The single-gluon normalisation of the setup (Sec.~Normalization of the main file) contains $\braket{(\Psi)_{gg}|(\Psi)_{gg}}=\frac{g^{2}N_{c}}{8\pi^{3}}\big(2L-\frac{11}{6}\big)\int\frac{d^{2}\tilde p}{\tilde p^{2}}$, set to zero as scaleless. In the cross section it enters as $(\mf{N}^{2}-1)\times$LO:
\begin{align}
\frac{d^{3}N_{\rm norm}}{d^{3}k}=-\frac{g^{2}N_{c}}{8\pi^{3}}\Big(2L-\frac{11}{6}\Big)\int\frac{d^{d}\tilde p}{\tilde p^{2}}\;\frac{d^{3}N_{\rm LO}}{d^{3}k}=\Big[\frac{\alpha_{s}\beta_{0}}{4\pi}-\mathfrak{a}L\Big]\Big(\frac{2}{\epsilon_{\rm UV}}-\frac{2}{\epsilon_{\rm IR}}\Big)\frac{d^{3}N_{\rm LO}}{d^{3}k}\qquad(=0\ \text{for}\ \epsilon_{\rm UV}=\epsilon_{\rm IR})\,.\label{eq:dv-norm}
\end{align}
\noindent[Calculation: $\frac{g^{2}N_{c}}{8\pi^{3}}\cdot\pi\cdot\frac{11}{6}=\frac{11g^{2}N_{c}}{48\pi^{2}}=\frac{\alpha_{s}\beta_{0}}{4\pi}$ and $\frac{g^{2}N_{c}}{8\pi^{3}}\cdot\pi\cdot2L=\frac{g^{2}N_{c}}{4\pi^{2}}L=\mathfrak{a}L$.] Its UV half, $\frac{\alpha_{s}N_{c}}{4\pi}\big(\frac{11}{3}-4L\big)\frac{2}{\epsilon_{\rm UV}}$, is the light-cone-gauge $Z_{3}-1$ of the produced gluon; its IR half ($\tilde p\to0$, collinear pair) is the virtual partner of the real splitting of Group II Row IV.

\subsection*{Step 4: Group II Row IV (real $g\to gg$ after the shock wave)}
\noindent The integral over the daughter separation $r=w-z$ at fixed parent positions diverges only at large $r$ (\texttt{s2}: $\pi\log\frac{R^{2}}{u^{2}}$, independent of the short-distance cut to $10^{-5}$): \textbf{IR (collinear)}; in $d=2-\epsilon$ it exists for $\epsilon<0$. Its pole multiplies the real splitting function $\mathcal{P}_{gg+}(\xi)=2N_{c}\big[\xi(1-\xi)+\frac{\xi}{1-\xi}+\frac{1-\xi}{\xi}\big]$ (no $\delta(1-\xi)$: a real emission). After the extraction of the end points the poles are: JIMWLK term $-(L+\ell)\,\mathfrak{a}\frac{2}{\epsilon}\frac{d^{3}N_{\rm LO}}{d^{3}k}$; DGLAP term
\begin{align}
-\mathcal{C}\frac{2}{\epsilon}\int_{0}^{1}d\xi\,\mathcal{P}_{gg+}\,\Pi=-\mathcal{C}\frac{2}{\epsilon}\big[\mathcal{P}_{gg}\otimes\Pi\big]+\frac{\alpha_{s}\beta_{0}}{4\pi}\frac{2}{\epsilon}\,\frac{d^{3}N_{\rm LO}}{d^{3}k}\,,\qquad\mathcal{P}_{gg}=\mathcal{P}_{gg+}+\frac{\beta_{0}}{2}\delta(1-\xi)\,,\qquad\mathcal{C}\frac{\beta_{0}}{2}\Pi(1)=\frac{\alpha_{s}\beta_{0}}{4\pi}\mathcal{N}_{0}\widehat{L}(k)\,,\label{eq:dv-IV}
\end{align}
the second term being the pole of the author's ``RC'' term: \textbf{its pole is collinear}, as the reviewer says. The $\xi\to0$ (BFKL) subtraction has no pole: the pole coefficient $\Pi(\xi)$ vanishes at $\xi=0$ for $k\neq0$.

\subsection*{Step 5: counterterms}
\noindent \emph{Coupling:} $g_{0}^{2}=\mu^{\epsilon}g^{2}(\mu)\big[1-\frac{\alpha_{s}\beta_{0}}{4\pi}\frac{2}{\epsilon}\big]$ in LO gives $-\frac{\alpha_{s}\beta_{0}}{4\pi}\frac{2}{\epsilon_{\rm UV}}\times$LO ($\rho$ and $U$ are not renormalised: the source is classical and $U$ is the given target field). \emph{Fragmentation} ($\overline{\rm MS}$, full $\mathcal{P}_{gg}$): $+\mathcal{C}\frac{2}{\epsilon_{\rm IR}}\big[\mathcal{P}_{gg}\otimes\Pi\big]$, with the $d$-dimensional Jacobian $\xi^{-d}$ (number-density FF; see the scheme note at the DGLAP term).

\subsection*{Step 6: the table}
\noindent Coefficients in units of $\mathfrak{a}\,\frac{2}{\epsilon}\times\frac{d^{3}N_{\rm LO}}{d^{3}k}$ (the units of the earlier note, $\frac{1}{(2\pi)^{3}}\frac{g^{4}N_{c}}{8\pi^{4}}\frac{1}{k^{+}}\frac{2}{\epsilon}\int e^{-ik\cdot(w-w')}\mathcal{O}(w,w')$, are the same), resp.\ of $\mathcal{C}\frac{2}{\epsilon}[\,\cdot\otimes\Pi]$ for the convolution terms; c.c.\ included for Group I.
\begin{center}
\renewcommand{\arraystretch}{1.3}\small
\begin{tabular}{|p{4.3cm}|p{3.6cm}|p{0.9cm}|p{2.9cm}|p{5.4cm}|}\hline
pole & origin (region) & UV/IR & coefficient & cancelled by / absorbed into\\\hline
Group I Row II, $\mf{A}^{(3)}$ row & $gg$ loop, $\tilde p\to\infty$ & UV & $\frac{11}{6}-L+\frac{\ell}{2}$ & $\frac{11}{6}$: Row V (Eqs.~\eqref{eq:dv-II}, \eqref{eq:dv-V}); $-L$: JIMWLK sector\\
Group I Row V, $2b$ UV & pair at one point, $z\to x$ & UV & $-\frac{11}{6}+2L$ & $-\frac{11}{6}$: Row II; $2L$: JIMWLK sector\\
dressed-gluon norm., UV half & $\tilde p\to\infty$ & UV & $\frac{11}{12}-L$ & coupling counterterm ($\frac{11}{12}$); $-L$: JIMWLK sector\\
coupling counterterm & $g_{0}\to g(\mu)$ & UV & $-\frac{11}{12}$ & UV half of the normalisation\\\hline
dressed-gluon norm., IR half & collinear pair, $\tilde p\to0$ & IR & $-\frac{11}{12}+L$ & $-\frac{11}{12}$: ``RC'' pole of Row IV; $L$: Row IV $\xi\to1$\\
Row IV, $\xi\to1$ end (JIMWLK term) & large $r$ & IR & $-L-\ell$ & $-L$: normalisation IR half; $-\ell$: left\\
Row IV, real $\mathcal{P}_{gg+}$ & large $r$ & IR & $-[\mathcal{P}_{gg}\otimes\Pi]+\frac{11}{12}$ & $-[\mathcal{P}_{gg}\otimes\Pi]$: FF counterterm; $+\frac{11}{12}$: normalisation IR half\\
Row IV, $\xi\to0$ (BFKL) & --- & --- & $0$ & ($\Pi(0)=0$ for $k\neq0$)\\
FF counterterm, full $\mathcal{P}_{gg}$ & $D_{g}\to D_{g}(\mu)$ & IR & $+[\mathcal{P}_{gg}\otimes\Pi]$ & real $\mathcal{P}_{gg+}$ $+$ virtual $\frac{\beta_{0}}{2}\delta$\\\hline
sum & & & UV: $\frac{\ell}{2}$; IR: $-\ell$ & $\to0$ for $\vee\to\infty$ (taken first)\\\hline
\end{tabular}
\end{center}
\noindent\textbf{Single-$\epsilon$ bookkeeping} (normalisation $=0$, the author's rule): Row II $+$ Row V $+$ Row IV $+$ coupling CT $+$ FF CT
\begin{align}
=\Big({\color{green!50!black}\cancel{\tfrac{11}{6}}}{\color{green!50!black}\cancel{-L}}+\tfrac{\ell}{2}\Big)+\Big({\color{green!50!black}\cancel{-\tfrac{11}{6}}}{\color{green!50!black}\cancel{+2L}}\Big)+\Big({\color{green!50!black}\cancel{-L}}-\ell{\color{green!50!black}\cancel{+\tfrac{11}{12}}}{\color{green!50!black}\cancel{-[\mathcal{P}_{gg}\otimes\Pi]}}\Big){\color{green!50!black}\cancel{-\tfrac{11}{12}}}{\color{green!50!black}\cancel{+[\mathcal{P}_{gg}\otimes\Pi]}}=-\frac{\ell}{2}\,,\notag
\end{align}
the same total as in the two-$\epsilon$ bookkeeping (UV $\frac{\ell}{2}$ $+$ IR $-\ell$). The $L$ poles cancel in the JIMWLK sector in both bookkeepings (UV: $-L+2L-L$; IR: $-L+L$). The remainder $-\frac{\ell}{2}\frac{2}{\epsilon}$ is a pole times a coefficient that vanishes for $\vee\to\infty$: it is zero because $\vee\to\infty$ is taken before $\epsilon\to0$ [\textbf{Answered} --- author's answer (date 2026-11-20): $\vee\to\infty$ before $\epsilon\to0$]. This order is required: the UV pole $+\frac{\ell}{2}\frac{2}{\epsilon}$ of Group I Row II cancels only for $\ell\to0$ first, since neither the range-restricted fragmentation counterterm nor the coupling counterterm supplies an $\ell$ term (\texttt{notes/scripts/s7\_range\_ct.py}, \texttt{s8\_range\_norm.py}). In the single-$\epsilon$ bookkeeping the coupling counterterm cancels the RC pole formally; this is legitimate only because the normalisation integral is scaleless (UV $=$ IR).

\subsection*{Step 7: physical account (answers to Q2 of the review)}
\begin{enumerate}
\item The virtual $\frac{11}{3}$ term of (I.89) is not missing: it is Eq.~\eqref{eq:dv-II}, cancelled exactly by Eq.~\eqref{eq:dv-V} (UV against UV). No finite term $\frac{\alpha_{s}\beta_{0}}{4\pi}[\mathcal{B}(x-w)+\mathcal{B}(x'-w')-\mathcal{B}(w-w')]\times$LO exists in this observable: the review is wrong on this point.
\item The review is right that the RC pole, as it arises in Group II Row IV, is collinear (Step 4). What renormalises the coupling is the UV half of the dressed-gluon normalisation, Eq.~\eqref{eq:dv-norm}, i.e.\ $Z_{3}$ of the produced gluon, consistent with the light-cone-gauge Ward identity $Z_{g}=Z_{3}^{-1/2}$ ($g_{0}A_{0}=gA$: the coupling of the classical source to the field needs no counterterm of its own). After the UV--UV cancellation of Rows II and V the only UV pole left is that of $Z_{3}$, $+\frac{\alpha_{s}\beta_{0}}{4\pi}\frac{2}{\epsilon_{\rm UV}}$ (the $L$ part belongs to the JIMWLK sector), which is exactly what $g_{0}^{2}\to g^{2}(\mu)$ removes.
\item The IR half of the same normalisation is the virtual $\frac{\beta_{0}}{2}\delta(1-\xi)$ (and $-2N_{c}L\,\delta(1-\xi)$) of the final-state splitting: real $\mathcal{P}_{gg+}$ $+$ virtual $=$ full $\overline{\rm MS}$ $\mathcal{P}_{gg}$, absorbed by the fragmentation counterterm. No IR pole is left once the FF counterterm with the full $\mathcal{P}_{gg}$ is subtracted.
\item ``Running coupling'': $N_{\rm RC}=-N_{\rm DGLAP}\big|_{\delta(1-\xi)}$ identically (added and subtracted). Its content is the $\mu$ dependence: $\partial_{\log\mu^{2}}N_{\rm RC}=+\frac{\alpha_{s}\beta_{0}}{4\pi}\times$LO compensates the running of $g^{2}(\mu)$ in LO, and the $\delta$ part of $N_{\rm DGLAP}$ compensates the $\delta$ part of the evolution of $D_{g}(\mu)$. In the classification of Kovner--Lublinsky--Skokov--Zhao (2308.15545, ``not all that is $\beta_{0}$ is $\beta$-function''): one $\beta_{0}$ (UV half of $Z_{3}$) is the $\beta$ function, the other (IR half, the $\frac{\beta_{0}}{2}\delta(1-\xi)$ of $\mathcal{P}_{gg}$) is DGLAP; here both are tied to the produced gluon and equal in magnitude because the integral is scaleless. In the local form (note at the running coupling section) LO $+$ RC $=$ LO $\times\frac{\alpha_{s}(k^{2})}{\alpha_{s}(\mu^{2})}$ to this order, plus a $\mu$-independent non-local remainder.
\end{enumerate}
}\par
